\documentclass[aspectratio=169,professionalfonts]{beamer}

\usetheme{Madrid}
\usecolortheme{dolphin}
\setbeamertemplate{navigation symbols}{}
\setbeamertemplate{footline}[frame number]

\usepackage[T1]{fontenc}
\usepackage[utf8]{inputenc}
\usepackage{lmodern}
\usepackage{booktabs}
\usepackage{array}
\usepackage{ragged2e}
\usepackage{xcolor}
\usepackage{graphicx}
\graphicspath{{docs/paper-review-skill-10min-talk/}{./}}

\definecolor{llnlblue}{RGB}{0,71,133}
\definecolor{llnlgold}{RGB}{235,170,0}
\definecolor{slate}{RGB}{45,55,72}
\definecolor{softgray}{RGB}{245,247,250}

\setbeamercolor{structure}{fg=llnlblue}
\setbeamercolor{title}{fg=white,bg=llnlblue}
\setbeamercolor{frametitle}{fg=white,bg=llnlblue}
\setbeamercolor{block title}{fg=white,bg=llnlblue}
\setbeamercolor{block body}{fg=slate,bg=softgray}
\setbeamercolor{alerted text}{fg=llnlgold}
\setbeamerfont{title}{series=\bfseries,size=\LARGE}
\setbeamerfont{frametitle}{series=\bfseries,size=\large}
\setbeamerfont{block title}{series=\bfseries}
\setbeamerfont{normal text}{size=\large}

\newcommand{\artifact}[1]{\texttt{\footnotesize #1}}
\newcommand{\claim}[1]{\textbf{\color{llnlblue}#1}}

\title[Paper Review Skill]{Paper Review Skill}
\subtitle{Auditable AI-assisted paper review with the \texttt{research-paper-review} Codex skill}
\author{Project talk draft}
\institute{paper-review-skill}
\date{\today}

\begin{document}

\begin{frame}
  \titlepage
  \note{
  Timing: about 45 seconds.

  Opening: this project is not just a prompt for reviewing papers. It is a
  local workflow that tries to make AI-assisted peer review inspectable,
  repeatable, and bounded by human reviewer responsibility.

  [Sources] README.md; SKILL.md.
  }
\end{frame}

\begin{frame}{The motivation was reproduction, then auditability}
  \begin{block}{Reference system}
    Biswas et al., \textit{AI-Assisted Peer Review at Scale: The AAAI-26 AI
    Review Pilot} (arXiv:2604.13940) describes an AI-assisted workflow for
    producing structured peer-review feedback at conference scale.
  \end{block}

  \vspace{0.7em}
  \begin{columns}[T,onlytextwidth]
    \column{0.47\textwidth}
      \claim{Reproduce}\\
      Implement the paper's staged review workflow as a Codex skill.
    \column{0.47\textwidth}
      \claim{Extend}\\
      Add OCR evidence, timing and provenance trails, rendered HTML, and a live
      web explainer.
  \end{columns}

  \note{
  Timing: about 1 minute.

  State the origin clearly: the first goal was to reproduce the workflow
  described in the AAAI-26 AI Review Pilot paper. The project then adds local
  engineering around that reproduced system: OCR evidence, timing and provenance
  logs, HTML rendering, and a live explainer. This sets up the next slide,
  which names the concrete artifact package.

  [Sources] README.md citation and project description; docs/audit-trails.md;
  docs/preprocessing.md.
  }
\end{frame}

\begin{frame}{The project turns one paper into an auditable review package}
  \begin{block}{One-sentence description}
    Paper Review Skill converts a paper into OCR-backed evidence, staged
    critique, timing and provenance logs, a final review, rendered HTML, and a
    local follow-up explainer.
  \end{block}

  \vspace{0.8em}
  \begin{tabular}{@{}p{0.28\textwidth}p{0.64\textwidth}@{}}
    \toprule
    Input & PDF, URL, Markdown, text, LaTeX, or review-form material \\
    Output & Review artifacts under \artifact{review\_artifacts/<paper\_id>/} \\
    Default backend & Authenticated \artifact{codex exec}; no separate API key required \\
    Completion bar & HTML report plus running local explainer server \\
    \bottomrule
  \end{tabular}

  \note{
  Timing: about 1 minute.

  This slide names the product boundary. The default workflow is built around
  Codex authentication and local scripts, not separate provider subscriptions.
  It supports URL inputs by downloading supported PDF URLs into the artifact
  tree before review.

  [Sources] README.md What It Does, Prerequisites, Basic Use, Main Outputs;
  docs/preprocessing.md.
  }
\end{frame}

\begin{frame}{The system separates inputs, processing, and outputs}
  \centering
  \includegraphics[width=0.98\textwidth,height=0.86\textheight,keepaspectratio]{paper-review-skill-system-diagram.pdf}

  \note{
  Timing: about 45 seconds.

  Use this slide to orient the audience before walking through the workflow.
  Inputs enter on the left, processing units run through the middle, and outputs
  leave on the right. Green parallelograms are inputs, blue rounded rectangles
  are processing units, gold folders are outputs, and the dashed gray node shows
  the shared Codex backend used by OCR, review drafting, and explainer Q\&A.

  [Sources] README.md Main Outputs; SKILL.md Default Workflow;
  docs/preprocessing.md; docs/audit-trails.md.
  }
\end{frame}

\begin{frame}{The workflow is staged so each claim has a trail}
  \begin{enumerate}
    \item Run mandatory OCR and HTML explainer preflights.
    \item Fetch or scope the PDF and produce canonical \texttt{olmOCR} Markdown.
    \item Draft focused stages: story, presentation, evaluation, correctness,
      and significance.
    \item Synthesize an initial review, audit it with self-critique, then revise.
    \item Run an independent quality critic before final rendering.
    \item Render HTML with artifacts exposed and start the explainer server.
  \end{enumerate}

  \note{
  Timing: about 1 minute.

  Present this as a pipeline of accountability rather than a chain of prompts.
  Each stage has a different review job, and the downstream artifacts make it
  possible to inspect where final comments came from.

  [Sources] SKILL.md Default Workflow; docs/current-skill-baseline.md Current
  Skill Workflow.
  }
\end{frame}

\begin{frame}{The design forces the boring parts to happen every time}
  \begin{columns}[T,onlytextwidth]
    \column{0.48\textwidth}
      \begin{block}{Mandatory checks}
        \begin{itemize}
          \item \artifact{check\_olmocr\_required.sh}
          \item \artifact{check\_html\_explainer\_required.sh}
          \item Stop on failure; do not fall back to static or PDF-only review
        \end{itemize}
      \end{block}
    \column{0.48\textwidth}
      \begin{block}{Default backend policy}
        \begin{itemize}
          \item \artifact{codex exec} for OCR shim, review drafting, and Q\&A
          \item \artifact{olmocr\_backend: codex-exec-shim}
          \item Hosted OCR providers are explicit overrides
        \end{itemize}
      \end{block}
  \end{columns}

  \vspace{0.5em}
  \begin{block}{First-principles reason}
    A trustworthy review aid must first preserve what was read, how it was
    transformed, and what system produced each judgment.
  \end{block}

  \note{
  Timing: about 1 minute.

  This is the core engineering choice: make dependency health, OCR evidence,
  and backend labeling non-optional. The repo instructions specifically call
  out that the default OCR backend should be recorded as codex-exec-shim, not
  as AllenAI olmOCR-2-7B-1025.

  [Sources] AGENTS.md project instructions; docs/preprocessing.md Mandatory
  Preflights and Important backend label.
  }
\end{frame}

\begin{frame}{The artifact tree is the audit interface}
  \begin{columns}[T,onlytextwidth]
    \column{0.52\textwidth}
      \begin{itemize}
        \item \artifact{evidence\_manifest.json}
        \item \artifact{model\_provenance.json}
        \item \artifact{ocr/<paper\_id>\_olmocr.md}
        \item \artifact{stages/*.md}
        \item \artifact{initial\_review.md}
        \item \artifact{self\_critique.md}
        \item \artifact{final\_review.md}
      \end{itemize}
    \column{0.43\textwidth}
      \begin{itemize}
        \item \artifact{quality\_report.md}
        \item \artifact{timing/timing.jsonl}
        \item \artifact{timing/timing\_summary.json}
        \item \artifact{timing/timing\_report.md}
        \item \artifact{timing/olmocr-pages.jsonl}
        \item \artifact{<paper\_id>\_review\_comments.html}
      \end{itemize}
  \end{columns}

  \vspace{0.4em}
  \begin{block}{Practical effect}
    The deliverable is not a single review blob; it is a review plus the
    evidence needed to challenge, revise, or reproduce it.
  \end{block}

  \note{
  Timing: about 1 minute.

  Explain that this directory structure gives both humans and tools stable
  places to inspect the workflow. It also supports resume behavior and batch
  processing because completion can be checked artifact by artifact.

  [Sources] README.md Main Outputs; docs/audit-trails.md Required Artifact
  Layout; docs/batch-processing.md Resume Behavior.
  }
\end{frame}

\begin{frame}{Review quality is treated as a separate artifact}
  \begin{columns}[T,onlytextwidth]
    \column{0.48\textwidth}
      \begin{block}{Stage coverage}
        \begin{itemize}
          \item End-to-end mechanism trace
          \item Technical soundness
          \item Evaluation and statistics
          \item Significance and positioning
          \item Writing and presentation
        \end{itemize}
      \end{block}
    \column{0.48\textwidth}
      \begin{block}{Quality gate}
        \begin{itemize}
          \item Missing required sections
          \item Unsupported criticism
          \item Identity leakage
          \item Rating/critique mismatch
          \item Broken Markdown or HTML
        \end{itemize}
      \end{block}
  \end{columns}

  \note{
  Timing: about 1 minute.

  The quality critic is not another paper reviewer. It reviews the generated
  review for completeness, grounding, appropriateness, and renderability. This
  separation keeps the pipeline from silently submitting a polished but weak
  or unsupported review.

  [Sources] SKILL.md Review Requirements; docs/quality-critic.md.
  }
\end{frame}

\begin{frame}{The HTML explainer closes the loop after rendering}
  \begin{block}{Why it matters}
    Reviewers rarely finish their thinking when the first draft appears. The
    local explainer lets them select review text or rendered artifacts and ask
    follow-up questions against the review package.
  \end{block}

  \vspace{0.7em}
  \begin{columns}[T,onlytextwidth]
    \column{0.33\textwidth}
      \claim{Render}\\
      HTML includes final review, staged artifacts, timing, and provenance.
    \column{0.33\textwidth}
      \claim{Serve}\\
      \artifact{start\_html\_explainer.sh} starts the local Q\&A endpoint.
    \column{0.33\textwidth}
      \claim{Record}\\
      Server status and URL are written back into the audit trail.
  \end{columns}

  \note{
  Timing: about 1 minute.

  The explainer is mandatory because the workflow expects review interrogation,
  not one-way text generation. The default URL is normally localhost on port
  8765, and batch mode can share one server across papers.

  [Sources] SKILL.md Interactive HTML Requirement; docs/preprocessing.md
  Evidence Manifest; docs/batch-processing.md Shared Explainer Server.
  }
\end{frame}

\begin{frame}{Engineering work keeps the review workflow maintainable}
  \begin{columns}[T,onlytextwidth]
    \column{0.5\textwidth}
      \begin{block}{Development discipline}
        \begin{itemize}
          \item Public repo excludes paper PDFs and generated private reviews
          \item Changes land through worktrees, PRs, CI, and regression checks
          \item Synthetic fixtures avoid committing real submissions
        \end{itemize}
      \end{block}
    \column{0.45\textwidth}
      \begin{block}{Validation surface}
        \begin{itemize}
          \item Skill eval validation
          \item Renderer regression fixtures
          \item Unit tests with coverage gate
          \item Script smoke tests
          \item Private eval runner for local papers
        \end{itemize}
      \end{block}
  \end{columns}

  \note{
  Timing: about 1 minute.

  This slide explains why the repository contains many small scripts and tests:
  the skill is an operational system. Public fixtures stay synthetic, while
  private evals can run locally against real papers without committing them.

  [Sources] docs/current-skill-baseline.md Repository State Policy, Validation
  and CI; docs/development-workflow.md; docs/private-evals.md.
  }
\end{frame}

\begin{frame}{What to remember}
  \begin{block}{Main takeaway}
    The project turns AI-assisted reviewing from a single model response into
    a reproducible, inspectable workflow with mandatory evidence, provenance,
    quality gates, and reviewer follow-up.
  \end{block}

  \vspace{0.8em}
  \begin{columns}[T,onlytextwidth]
    \column{0.31\textwidth}
      \claim{For reviewers}\\
      Faster synthesis, but every judgment remains theirs.
    \column{0.31\textwidth}
      \claim{For maintainers}\\
      Stable artifacts make regressions visible.
    \column{0.31\textwidth}
      \claim{For deployment}\\
      Policy, confidentiality, and venue rules still decide when to use it.
  \end{columns}

  \note{
  Timing: about 45 seconds, leaving time for one or two questions.

  Close by returning to responsibility. The skill can help structure review
  work, but it does not make venue policy or reviewer accountability disappear.
  A good next discussion is which private evals or real review forms should be
  used to test whether the workflow is actually useful.

  [Sources] README.md Responsible Use; docs/private-evals.md.
  }
\end{frame}

\end{document}
